\subsection{Linear complex structures}

\bdefn

    A {\emphcolor (linear) complex structure} on a real vector space is an automorphism
    $$ J\colon V\to V \qquad\hbox{such that}\qquad J^2 = -I . $$
    Denote the space of linear complex structures over $V$ by
    $$ \J(V) \subseteq \endo(V) . $$

\edefn

The archetypal linear complex structure is
$$ J_0 = \pmatrix{0&-I\cr I&0} $$
over $\bR^{2n}$.

$V$ having a complex structure is equivalent to it having the structure of a $\bC$-vector space compatible with its real vector space structure.
This is equivalent to the existence of a ring map $\bC\to\endo(V)$ which is a lifting of the ring map $\bR\to\endo(V)$ which defines $V$'s real linear structure:

\diagram{
    \bR&\endo(V)\cr
    \bC
}{
    \diagarrow{from={1,1}, to={1,2}}
    \diagarrow{from={1,1}, to={2,1}, left cap={(}}
    \diagarrow{from={2,1}, to={1,2}}
}

That is, the existence of a map of $\bR$-algebras $\bC\to\endo(V)$.
Because $\bC$ is generated by $i$ as a real algebra, this is equivalent to the existence of $J\in\endo(V)$ such that $J^2=-I$.
Spelling this out, if $V$ has a complex structure then it becomes a complex vector space with multiplication
$$ \bC\times V\to V,\qquad (a+ib,v) \mapsto av + bJv . $$

Notice that this necessitates that $V$ has even real dimension.
This can be seen because
$$ \dim_\bR V = \dim_\bC V\cdot\dim_\bR\bC = 2\dim_\bC V $$
and because $J^2=-I$ means that $\det J^2=(-1)^{\dim V}$ is positive so $\dim V$ is even.

\bprop

    Let $V$ be a $2n$-dimensional real vector space with complex structure $J$.
    Then there is a vector space isomorphism $\Phi\colon\bR^{2n}\to V$ such that $J\Phi=\Phi J_0$.

\eprop

\bproof

    Viewing $V$ as a complex vector space, there is a basis $v_1,\dots,v_n$.
    Then $v_1,\dots,v_n,Jv_1,\dots,Jv_n$ is a basis of the real vector space $V$.
    Mapping the standard basis $e_1,\dots,e_n,f_1,\dots,f_n$ to this basis gives the desired result.
    \qed

\eproof

